En este informe se estudió experimentalmente la ley de enfriamiento de Newton mediante el análisis de la variación de temperatura de un sistema en función del tiempo. Para ello, se calentaron \(350\ \mathrm{mL}\) de agua y posteriormente se registró su proceso de enfriamiento hasta aproximarse a la temperatura ambiente. Los datos obtenidos fueron representados gráficamente como temperatura contra tiempo y también mediante la linealización de \(\ln(T-T_a)\) en función del tiempo.

Los resultados mostraron un comportamiento exponencial consistente con el modelo teórico de la ley de enfriamiento de Newton. La gráfica linealizada presentó un ajuste satisfactorio, con un coeficiente de correlación \(r^2 = 0{,}99998230\), lo cual evidencia una fuerte concordancia entre la teoría y los datos experimentales. Finalmente, se analizaron las posibles fuentes de incertidumbre asociadas al experimento, concluyendo que estas no afectaron significativamente la tendencia general observada.
